\section{Stationary states}

\SourcePage{44}{47}
The Hamiltonian of a closed system (and also of a system in a constant, but not
a variable, external field) cannot contain time explicitly. This follows
because all instants of time are equivalent for such a physical system. Since,
on the other hand, every operator commutes with itself, we conclude that the
Hamilton function is conserved for systems not situated in a variable external
field. A conserved Hamilton function is called the energy. The meaning of
energy conservation in quantum mechanics is that if the energy has a definite
value in a given state, this value remains constant in time.

States possessing definite energy values are called the \textit{stationary
states} of the system. They are described by wave functions $\Psi_n$ that are
eigenfunctions of the Hamiltonian operator, satisfying
$\hat H\Psi_n=E_n\Psi_n$, where $E_n$ are the energy eigenvalues. Accordingly,
the wave equation (8.1) for
\SourcePage{45}{48}
the function $\Psi_n$,
\[
  i\hbar\frac{\partial\Psi_n}{\partial t}=\hat H\Psi_n=E_n\Psi_n,
\]
can be integrated directly with respect to time, giving
\begin{equation}
  \Psi_n=\exp\left(-\frac{i}{\hbar}E_nt\right)\psi_n(q),
  \tag{10.1}
\end{equation}
where $\psi_n$ is a function of the coordinates alone. This determines the
time dependence of the wave functions of stationary states.

We shall use the lower-case letter $\psi$ for stationary-state wave functions
without the time factor. These functions and the energy eigenvalues themselves
are determined by
\begin{equation}
  \hat H\psi=E\psi.
  \tag{10.2}
\end{equation}
Among all stationary states, the one possessing the lowest energy is termed
the \textit{normal}, or \textit{ground}, state of the system.

The expansion of any wave function $\Psi$ over the stationary-state wave
functions has the form
\begin{equation}
  \Psi=\sum_n a_n\exp\left(-\frac{i}{\hbar}E_n(t)\right)\psi_n(q).
  \tag{10.3}
\end{equation}
As usual, the squares $|a_n|^2$ give the probabilities of the various energy
values of the system.

The coordinate probability distribution in a stationary state is determined
by $|\Psi_n|^2=|\psi_n|^2$ and is therefore independent of time. The same is
true of the mean value
$\bar f=\int\Psi_n^*\hat f\Psi_n,dq=\int\psi_n^*\hat f\psi_n,dq$ of any
physical quantity $f$ whose operator does not depend explicitly on time.

As stated above, the operator of every conserved quantity commutes with the
Hamiltonian. Thus every conserved physical quantity can be measured
simultaneously with the energy.

Among the stationary states there may be states corresponding to one and the
same energy eigenvalue (or \textit{energy level}) but differing in the values
of other physical quantities. A level to which several distinct stationary
states correspond is said to be \textit{degenerate}. Physically, degeneracy is
possible because the energy does not, in general, constitute by itself a
complete set of physical quantities.
\SourcePage{46}{49}
The energy levels of a system are, in general, degenerate if there are two
conserved quantities $f$ and $g$ whose operators do not commute. Let $\psi$ be
the wave function of a stationary state in which $f$, as well as the energy,
has a definite value. Then $\hat g\psi$ cannot coincide with $\psi$ up to a
constant factor; otherwise $g$ too would have a definite value, which is
impossible because $f$ and $g$ cannot be measured simultaneously. On the other
hand, $\hat g\psi$ is a Hamiltonian eigenfunction with the same energy $E$ as
$\psi$:
\[
  \hat H(\hat g\psi)=\hat g\hat H\psi=E(\hat g\psi).
\]
Thus more than one eigenfunction corresponds to $E$, and the level is
degenerate.

Any linear combination of wave functions belonging to the same degenerate
energy level is also an eigenfunction with that energy. The choice of
eigenfunctions for a degenerate eigenvalue is therefore not unique. Arbitrarily
chosen eigenfunctions of a degenerate level are not in general mutually
orthogonal, but suitable linear combinations can always be chosen to form a
mutually orthogonal and normalized set\footnote{This can be done in infinitely
many ways: a linear transformation of $n$ functions involves $n^2$ independent
coefficients, but the normalization and orthogonality conditions number only
$n(n+1)/2$, which is fewer than $n^2$.}.

These statements apply not only to energy eigenfunctions but to the
eigenfunctions of any operator. Only functions belonging to different
eigenvalues are automatically orthogonal; functions belonging to the same
degenerate eigenvalue are not in general orthogonal.

If $\hat H=\hat H_1+\hat H_2$, where one part contains only coordinates $q_1$
and the other only $q_2$, the eigenfunctions of $\hat H$ can be written as
products of eigenfunctions of $\hat H_1$ and $\hat H_2$, and the energy
eigenvalues are sums of the eigenvalues of these operators.
\SourcePage{47}{50}
The energy spectrum may be discrete or continuous. A stationary state in the
discrete spectrum always corresponds to \textit{finite motion}: neither the
system nor any part of it recedes to infinity. Indeed,
$\int|\Psi|^2,dq$ is finite over all space for a discrete-spectrum
eigenfunction. Thus $|\Psi|^2$ decreases sufficiently rapidly and vanishes at
infinity. Equivalently, the probability of infinite coordinate values is zero;
the system performs finite motion, or is in a \textit{bound} state.

For continuous-spectrum wave functions, $\int|\Psi|^2,dq$ diverges.
$|\Psi|^2$ then determines coordinate probabilities only up to
proportionality. The divergence is always associated with $|\Psi|^2$ failing
to vanish at infinity, or vanishing too slowly. The integral over the region
outside any arbitrarily large but finite closed surface therefore still
diverges: the system, or some part of it, is at infinity. A superposition of
continuous-spectrum stationary states may have a convergent integral and be
localized in a finite region, but this region spreads without bound with time,
and ultimately the system recedes to infinity.

Indeed, an arbitrary continuous-spectrum superposition is
\[
  \Psi=\int a_E\exp\left(-\frac{i}{\hbar}Et\right)\psi_E(q)\,dE,
\]
and
\[
  |\Psi|^2=\iint a_Ea_{E'}^*\exp\left(\frac{i}{\hbar}(E'-E)t\right)
  \psi_E(q)\psi_{E'}^*(q)\,dE\,dE'.
\]
If this expression is averaged over an interval $T$ and then $T$ is allowed to
tend to infinity, the oscillating factors average to zero, and with them
\SourcePage{48}{51}
the entire integral. Thus the time-averaged probability of finding the system
at any specified point of configuration space vanishes; this is possible only
if the motion occupies all infinite space\footnote{For a superposition of
discrete-spectrum functions one would instead have
\[
  \overline{|\Psi|^2}=\overline{\sum_{n,m}a_na_m^*
  \exp\left\{\frac{i}{\hbar}(E_m-E_n)t\right\}\psi_n\psi_m^*}
  =\sum_n|a_n\psi_n(q)|^2,
\]
so the probability density remains finite under time averaging.}.

Thus stationary states of the continuous spectrum correspond to infinite
motion of the system.
